\begin{comment}
\subsubsection{Bayesian Test}

One problem with the $G$ test is its finite-sample behaviour. 
One can ``smoothen'' it by using a related Bayesian test.
We will compare two Bayesian models for the joint distribution $P(X,Y)$.
In the ``dependent'' model we take the same likelihood as before, and now
we put a Dirichlet prior on the parameters $\theta_{xy} \sim \mathrm{Dir}(\alpha_{xy})$. 
We can think of the hyperparameters $\alpha_{xy}$ as ``pseudo-counts''.
The Bayesian test can be implemented as model comparison for two (nested) models, one corresponding with the
null hypothesis of independence, i.e., with $\theta_{xy} = \theta_{+ x} \theta_{y +}$, and one without
such a restriction on the parameters. 
The \emph{Bayes factor} is the ratio of the two \emph{model evidences} (\emph{marginal likelihoods}).
$$BF = \frac{\int P_{H_1}((X_n,Y_n)_{n=1}^N \mid \theta) dP_{H_1}(\theta)}{\int P_{H_0}((X_n,Y_n)_{n=1}^N \mid \theta) dP_{H_0}(\theta)}$$
The likelihood ratio is replaced by the ratio of the model evidences:

\begin{Def}
  The \emph{Dirichlet distribution} with parameters $\alpha_a,\dots,\alpha_K > 0$
  has probability density with respect to Lebesgue measure in $\RN^{K-1}$ given by
\end{Def}

Dirichlet prior:
$$\mathrm{Dir}(\B{\mu} \given \B{\alpha}) = \frac{\Gamma(\sum_k \alpha_k)}{\prod_k \Gamma(\alpha_k)} \prod_k \mu_k^{\alpha_k - 1}$$
Multinomial likelihood:
$$\mathrm{Mult}(\B{N} \given \B{\mu}) = \frac{N!}{\prod_k N_k!} \prod_k \mu_k^{N_k}$$
The marginal likelihood of the \emph{data} is:
$$p(\B{X} \given \B{\alpha}) = \frac{\Gamma(\sum_k \alpha_k)}{\prod_k \Gamma(\alpha_k)} \frac{\prod_k \Gamma(\alpha_k + N_k)}{\Gamma(N + \sum_k \alpha_k)}$$
The marginal likelihood of the \emph{counts} is:
$$p(\B{N} \given \B{\alpha}) = \int \mathrm{Mult}(\B{N} \given \B{\mu}) \mathrm{Dir}(\B{\mu} \given \B{\alpha}) d\B{\mu} = \frac{\Gamma(\sum_k \alpha_k)}{\prod_k \Gamma(\alpha_k)} \frac{\Gamma(N + 1)}{\prod_k \Gamma(N_k + 1)} \frac{\prod_k \Gamma(\alpha_k + N_k)}{\Gamma(N + \sum_k \alpha_k)}$$
For the special case that $\B{\alpha} = \B{1}$, this simplifies to
$$p(\B{N} \given \B{\alpha} = 1) = \frac{\Gamma(K) \Gamma(N + 1)}{\Gamma(N + K)}$$
where $K := \sum_k 1$.

Taking $(X,Y)$ to have a multinomial distribution with parameters $\mu^{XY} \sim \mathrm{Dir}(\B{\alpha}^{XY})$, we obtain the following marginal likelihood:
$$p(\B{N}^{XY} \given \B{\alpha}^{XY}) = \frac{\Gamma(\sum_k \sum_l \alpha^{XY}_{kl})}{\prod_k \prod_l \Gamma(\alpha^{XY}_{kl})} \frac{\Gamma(N^{XY} + 1)}{\prod_k \prod_l \Gamma(N^{XY}_{kl} + 1)} \frac{\prod_k \prod_l \Gamma(\alpha^{XY}_{k,l} + N^{XY}_{k,l})}{\Gamma(N^{XY} + \sum_k \sum_l \alpha^{XY}_{kl})}$$

Now taking $(X,Y)$ to have a multinomial distribution with parameters $(\mu^{XY})_{kl} = (\mu^X)_k (\mu^Y)_l$, where $\mu^X \sim \mathrm{Dir}(\B{\alpha}^X)$ and
$\mu^Y \sim \mathrm{Dir}(\B{\alpha}^Y)$, we obtain:
\begin{equation*}\begin{split}
& p(\B{N}^{XY} \given \B{\alpha}^X, \B{\alpha}^Y) \\
& = \int \frac{\Gamma(N^{XY} + 1)}{\prod_k \prod_l \Gamma(N^{XY}_{kl} + 1)} \prod_k \prod_l (\mu^X_k \mu^Y_l)^{N^{XY}_{kl}} 
  \frac{\Gamma(\sum_k \alpha^X_k)}{\prod_k \Gamma(\alpha^X_k)} \prod_k (\mu^X_k)^{\alpha^X_k - 1} \frac{\Gamma(\sum_l \alpha^Y_l)}{\prod_l \Gamma(\alpha^Y_l)} \prod_l (\mu^Y_l)^{\alpha^X_l - 1} d\B{\mu^X} d\B{\mu^Y} \\
& = \frac{\Gamma(N^{XY} + 1)}{\prod_k \prod_l \Gamma(N^{XY}_{kl} + 1)} \frac{\Gamma(\sum_k \alpha^X_k)}{\prod_k \Gamma(\alpha^X_k)} \frac{\Gamma(\sum_l \alpha^Y_l)}{\prod_l \Gamma(\alpha^Y_l)} 
  \int \prod_k \prod_l (\mu^X_k \mu^Y_l)^{N^{XY}_{kl}} \prod_k (\mu^X_k)^{\alpha^X_k - 1} \prod_l (\mu^Y_l)^{\alpha^X_l - 1} d\B{\mu^X} d\B{\mu^Y} \\
& = \frac{\Gamma(N^{XY} + 1)}{\prod_k \prod_l \Gamma(N^{XY}_{kl} + 1)} \frac{\Gamma(\sum_k \alpha^X_k)}{\prod_k \Gamma(\alpha^X_k)} \frac{\Gamma(\sum_l \alpha^Y_l)}{\prod_l \Gamma(\alpha^Y_l)}
  \int \prod_k (\mu^X_k)^{N^X_k + \alpha^X_k - 1} d\B{\mu^X} \int \prod_l (\mu^Y_l)^{N^Y_l + \alpha^X_l - 1} d\B{\mu^Y} \\
& = \frac{\Gamma(N + 1)}{\prod_k \prod_l \Gamma(N^{XY}_{kl} + 1)} \frac{\Gamma(\sum_k \alpha^X_k)}{\prod_k \Gamma(\alpha^X_k)} \frac{\Gamma(\sum_l \alpha^Y_l)}{\prod_l \Gamma(\alpha^Y_l)}
  \frac{\prod_k \Gamma(N^X_k + \alpha^X_k)}{\Gamma(N + \sum_k \alpha^X_k)} \frac{\prod_l \Gamma(N^Y_l + \alpha^Y_l)}{\Gamma(N + \sum_l \alpha^Y_l)} 
\end{split}\end{equation*}
where $N = \sum_k N^X_k = \sum_l N^Y_l = \sum_{k,l} N^{XY}_{kl}$.

So the Bayes factor (independence vs.\ dependence) is given by:
\begin{equation}\label{eq:Bayesian_indep_test}\begin{split}
& \frac{p(\B{N}^{XY} \given \B{\alpha}^X, \B{\alpha}^Y)}{p(\B{N}^{XY} \given \B{\alpha}^{XY})} = 
  \frac{\prod_k \prod_l \Gamma(\alpha^{XY}_{kl})}{\Gamma(\sum_k \sum_l \alpha^{XY}_{kl})} \frac{\Gamma(N + \sum_k \sum_l \alpha^{XY}_{kl})}{\prod_k \prod_l \Gamma(\alpha^{XY}_{k,l} + N^{XY}_{k,l})} \\
& \frac{\Gamma(\sum_k \alpha^X_k)}{\prod_k \Gamma(\alpha^X_k)} \frac{\Gamma(\sum_l \alpha^Y_l)}{\prod_l \Gamma(\alpha^Y_l)} \frac{\prod_k \Gamma(N^X_k + \alpha^X_k)}{\Gamma(N + \sum_k \alpha^X_k)} \frac{\prod_l \Gamma(N^Y_l + \alpha^Y_l)}{\Gamma(N + \sum_l \alpha^Y_l)} \\
\end{split}\end{equation}
This looks pretty complicated, but is straightforward to compute. It simplifies considerably if we take all $\alpha$ components to be 1:
\begin{equation*}\begin{split}
& \frac{p(\B{N}^{XY} \given \B{\alpha}^X = \B{1}, \B{\alpha}^Y = \B{1})}{p(\B{N}^{XY} \given \B{\alpha}^{XY} = \B{1})} = 
  \frac{\Gamma(|X|) \Gamma(|Y|)}{\Gamma(|X||Y|)} \frac{\Gamma(N + |X||Y|)}{\prod_k \prod_l \Gamma(N^{XY}_{k,l} + 1)} 
  \frac{\prod_k \Gamma(N^X_k + 1)}{\Gamma(N + |X|)} \frac{\prod_l \Gamma(N^Y_l + 1)}{\Gamma(N + |Y|)} \\
\end{split}\end{equation*}

\Joris{Approach: obtain asymptotic marginal likelihoods by BIC. It uses a similar regularity condition, hence we will again
only consider the case with positive parameters. The advantage is that the proof is pretty short, so we may be able to figure
out what changes in the ``conditional'' setting. A few steps todo: link the Dirichlet-multinomial to the notation and
assumptions in Schwartz 1978. We still have to solve the ML estimate, and then correct it with a BIC. This will turn out to
be a simple correction on the likelihood ratio, but with pseudocounts added. We will get the most general approach if we do
the asymptotic analysis ``by hand'', not relying too much on standard statistical results.

With the BIC correction, we get an additional
correction term. It says: log marginal likelihood converges to maximized log likelihood
minus half dim $\log N$ plus order one remainder. But that means that the log Bayes factor
converges to the log likelihood ratio plus some $O(\log N)$ correction. Converges in which sense?
Note that under $H_1$, $G_N = O_P(N)$, and hence $G_N \stackrel{P}{\to} \infty$, and therefore
$BF_N \stackrel{P}{\to} +\infty$, while under $H_0$, $G_N = O_P(1)$, and hence $BF_N \stackrel{P}{\to} -\infty$.
}
\end{comment}

\subsubsection{Marginal Independence of a Random and a Non-Random Variable}

Consider two variables $X,C$ taking values in finite spaces $\C{X},\C{C}$, respectively,
(i.e., with $2\le|\C{X}|<\infty$ and $2\le|\C{C}|<\infty$).
Assume that $X$ is a random variable, while $C$ is an input variable, 
and consider a Markov kernel $K(X\mid C) : \C{C} \dshto \C{X}$. 
We will derive a test for the independence $$X \Indep_{K(X\mid C)} C.$$
%We can represent this Markov kernel in a table, e.g., for the binary case:
%\begin{center}\begin{tabular}{l|ll}
%        & $C=0$         & $C=1$ \\
%        \hline
%  $X=0$ & $p_{0|0}$      & $p_{0|1}$ \\
%  $X=1$ & $p_{1|0}$      & $p_{1|1}$ \\
%  \hline
%        & 1              & 1
%\end{tabular}\end{center}
%where $p_{x|c} = K(X=x|C=c)$. 
With Definition~\ref{cond-indep-markov-kernels-def}, this independence holds
%$$X \Indep_{K(X|C)} C$$
if and only if there exists a Markov kernel $Q(X) : \Asterisk \dshto \C{X}$ such that:
\[ K(X\mid C) = Q(X).\]
The latter means that
\begin{equation}\label{eq:indep_test_CX}
  K(X \mid C=c) = Q(X)
\end{equation}
for all $c \in \C{C}$.
%For a binary context, the independence is equivalent to $p_{x|0} = p_{x|1}$ for all $x \in \C{X}$.

Suppose we obtain data $(X_n,c_n)_{n=1}^N$ such that the
$X_n$ are conditionally independent and identically distributed given $c_n$ for all $n=1,\dots,N$.
In other words, we assume the data is sampled from the following Markov kernel:
$$K((X_n)_{n=1}^N \mid (c_n)_{n=1}^N) = \bigotimes_{n=1}^N K(X=X_n \mid C=c_n),$$
with the ``true'' kernel $K(X \mid C)$ unknown.

Note that this is a weaker assumption regarding the sampling scheme than
we would have made if $C$ were random. In particular, we make no assumption
at all regarding how the sequence of values $c_1,c_2,\dots,c_N$ is chosen.
It could be a sequence like $0,1,0,1,0,1,0,1,\dots$, for example, which would
be (if sufficiently long) very unlikely to occur if all $c_n$ would be independently
sampled from some distribution. This broadens the possible experimental designs that
we can handle to include for example randomized controlled trials in which the 
protocol is such that a certain prespecified number $N_{+|0}$ of subjects enters the control
group, and a certain prespecified number $N_{+|1}$ enters the treatment group. If the values
of $c_n$ were chosen i.i.d.\ with a coin flip, then it would be very unlikely that
this assignment satisfies the protocol. Considering $C$ to be an exogenous input
variable instead (with values that are not necessarily randomly assigned), allows us to test for the 
independence of outcome $X$ and treatment $C$ under a broader range of experimental
protocols.

We define the \emph{counts} as the number of observations with a given value $(x,c)\in\C{X}\times\C{C}$:
$$N_{x|c} := \sum_{n=1}^N \ind_{(x,c)}(X_n,c_n).$$
%In the binary case, we can represent them in a table:
%\begin{center}\begin{tabular}{l|ll}
%        & $C=0$         & $C=1$ \\
%        \hline
%  $X=0$ & $N_{0|0}$      & $N_{0|1}$     \\
%  $X=1$ & $N_{1|0}$      & $N_{1|1}$     \\
%  \hline
%        & $N_{+|0}$ & $N_{+|1}$ 
%\end{tabular}\end{center}
%where we used a similar summation convention as before:
%$$N_{+|c} := \sum_{x\in\C{X}} N_{x|c}.$$

We will take $H_0 : X \Indep C$ as the null hypothesis of a frequentist test for the independence of $X$ and $C$.
We parameterize the Markov kernel $K(X \mid C)$ in terms of parameters
$(\theta_{x|c})_{x\in\C{X},c\in\C{C}} := K(X = x \mid C = c)$ in a space
$$\Theta := %\{ \theta \in \prod_{z\in\C{Z}} [0,1] : \sum_{z\in\C{Z}} \theta_z = 1\} \times
\{ \theta \in \prod_{x\in\C{X},c\in\C{C}} [0,1] : \forall_{c\in\C{C}} \sum_{x\in\C{X}} \theta_{x|c} = 1\}.$$
With the summation convention, we can write the normalization condition as $\theta_{+|c} = 1$ for all $c \in \C{C}$.
The null hypothesis $H_0 : X \Indep C$, equivalent to \eref{eq:indep_test_CX}, can be expressed
in terms of the parameters as $H_0 : \theta \in \Theta_0$, where we introduced the restricted parameter space
$$\Theta_0 := \{ \theta \in \Theta : \forall_{x\in\C{X},c\in\C{C},c'\in\C{C}}: \theta_{x|c} = \theta_{x|c'} \}.$$
%which is also equivalent to
%$$\exists_{\theta_x} : \forall_{x\in\C{X},c\in\C{C}}: \theta_{x|c} = \theta_{x}.$$
As alternative hypothesis we will take $H_1 : X \nIndep C$, the negation of the null hypothesis, i.e.,
$H_1 : \theta \in \Theta_1$ with $\Theta_1 := \Theta \setminus \Theta_0$.

We will again work out the likelihood ratio test. 
We first write down the ``conditional'' likelihood of the data:
$$K\big((X_n)_{n=1}^N \mid (c_n)_{n=1}^N, \theta\big)
= \prod_{n=1}^N \theta_{X_n|c_n} 
= \prod_{x\in\C{X}} \prod_{c\in\C{C}} \theta_{x|c}^{N_{x|c}}$$
where we used the counts as a sufficient statistic of the data.
Maximizing the likelihood with respect to the parameters, we obtain the maximum likelihood estimator
$$\hat \theta_{x|c} = \frac{N_{x|c}}{N_{+|c}},$$
i.e., the fractions of the different outcomes within each subgroup with $C = c$.
Under $H_0$, we can write $\theta_{x|c} = \theta_{x|\ast}$ for some $\theta_{x|\ast} \in \Theta_0$, 
and the likelihood simplifies:
$$\theta \in \Theta_0 \implies \prod_{x\in\C{X}} \prod_{c\in\C{C}} \theta_{x|\ast}^{N_{x|c}} = 
\prod_{x\in\C{X}} \theta_{x|\ast}^{N_{x|+}}.$$
The maximum likelihood estimator under $H_0$ is then
$$\hat \theta^0_{x|c} = \frac{N_{x|+}}{N}.$$

The likelihood ratio is obtained by dividing the likelihood for $\hat \theta$ by the likelihood for $\hat \theta^0$:
\begin{equation}\label{eq:ext_G_test_statistic}
  \frac{\sup_{\theta\in\Theta_0 \cup \Theta_1} K((X_n)_{n=1}^N \mid (c_n)_{n=1}^N, \theta)}{\sup_{\theta\in\Theta_0} K((X_n)_{n=1}^N \mid (c_n)_{n=1}^N, \theta)}
= \prod_{x\in\C{X}} \prod_{c\in\C{C}} \left(\frac{\hat \theta_{x|c}}{\hat \theta^0_{x|c}}\right)^{N_{x|c}} 
= \prod_{x\in\C{X}} \prod_{c\in\C{C}} \left(\frac{N_{x|c}}{N_{+|c}} \frac{N}{N_{x|+}}\right)^{N_{x|c}}.
\end{equation}
This is of \emph{exactly} the same form as the likelihood ratio \eref{eq:indep_test_XY_likelihood_ratio} for testing
$X \Indep_{P(X,Y)} Y$ with two random variables $X,Y$.
The likelihood ratio test statistic is defined as 2 times the natural logarithm of this ratio:
$$G_N := 2 \log \frac{\sup_{\theta\in\Theta_0 \cup \Theta_1} K((X_n)_{n=1}^N \mid (c_n)_{n=1}^N, \theta)}{\sup_{\theta\in\Theta_0} K((X_n)_{n=1}^N \mid (c_n)_{n=1}^N, \theta)} = 2 \sum_{x\in\C{X}} \sum_{c\in\C{C}} N_{x|c} \log \frac{N_{x|c} N}{N_{x|+} N_{+|c}}$$

We (miraculously?!) arrived at the same test statistic as before.
This time, we cannot make use of the general result on the asymptotic distribution under the null hypothesis of 
likelihood ratio tests. Indeed, that result pertains when dealing
with a likelihood with only a finite number of parameters to be estimated, whereas here we have an asymptotically infinite number
of ``parameters'' $c_1,c_2,\dots$. Therefore, we will resort to a more direct analysis of the case at hand. The end result will
recover the previous results for random variables as a special case. Perhaps surprisingly, it turns out that under reasonable 
assumptions on the sequence $c_1,c_2,\dots$ we can apply the standard $G$ test and ignoring the non-random nature of the $c_n$'s.

We will start with rewriting the test statistic. We introduce the space
$$\Theta_C := \{ \gamma \in \prod_{c\in\C{C}} [0,1] : \sum_{c\in\C{C}} \gamma_c = 1 \}.$$
Consider now the function
\begin{equation}\label{eq:G_test_statistics_limit}
  \begin{split}
  g : \Theta \times \Theta_C : (\theta,\gamma) \mapsto
    &:= \sum_{c\in\C{C}} \gamma_c \sum_{x\in\C{X}} \theta_{x|c} \log \frac{\theta_{x|c}}{\sum_{c\in\C{C}} \gamma_c \theta_{x|c}} \\
    & = \sum_{c\in\C{C}} \gamma_c KL\Big(\theta_{X|c} \,\Vert\, \sum_{c\in\C{C}} \gamma_c \theta_{X|c}\Big),
\end{split}\end{equation}
where we introduced the notation $\theta_{X|c} := (\theta_{x|c})_{x\in\C{X}} \in \RN^{\C{X}}$, and where
the Kullback-Leibler divergence between two probability distributions $P,Q \in \C{P}(\C{X})$ is defined as:
$$KL(P \,\Vert\, Q) := \sum_{x\in\C{X}} P(X=x) \log \frac{P(X=x)}{Q(X=x)}$$
(and we identified a probability distribution on $\C{X}$ with its probability mass function, encoded as a parameter vector in 
$\RN^\C{X}$).
Note that
$$G_N = 2 N g(\hat\theta, \hat\gamma)$$
with
$$\hat\gamma_c := \frac{N_{+|c}}{N}$$
for all $c\in\C{C}$.\footnote{We write $\hat\gamma$ rather than $\gamma$ to indicate that it is an ($N$-dependent) function of the observed data, rather than a ``true'' fixed quantity.}
The components of $\hat\gamma$ are just the fractions of observations with a certain value of $c$. 
The asymptotic analysis will be conditional on the sequence of $\hat\gamma$'s, or equivalently,
on the sequence of counts $N_{+|c}$.

Because the counts $(N_{x|c})_{x\in\C{X}}$ have a multinomial distribution for each $c\in\C{C}$, we can calculate that
$$\Exp \hat\theta_{x|c} = \frac{1}{N_{+|c}} \sum_{n=1}^N \I_c(c_n) \theta_{x|c} = \theta_{x|c}$$
and
\begin{equation*}\begin{split}
  \Cov(\hat \theta_{x|c}, \hat \theta_{x'|c'}) 
  & = \frac{1}{N_{+|c} N_{+|c'}} \sum_{n=1}^N \sum_{n'=1}^N \Cov\big( \I_{(x,c)}(X_n,c_n) \I_{(x',c')}(X_{n'},c_{n'}) \big) \\
%  & = \frac{1}{N_{+|c} N_{+|c'}} \sum_{n=1}^N \sum_{n'=1}^N \I_c(c_n) \I_{c'}(c_{n'}) \big( P(X_n=x,X_{n'}=x' \mid c_n=c,c_{n'}=c') - P(X_n=x \mid c_n=c) P(X_{n'}=x' \mid c_{n'}=c') \big) \\
  & = \frac{1}{N_{+|c} N_{+|c'}} \sum_{n=1}^N \I_c(c_n) \I_{c'}(c_{n}) \big( \delta_{xx'} \theta_{x|c} - \theta_{x|c} \theta_{x'|c} \big) \\
  & = \frac{1}{N_{+|c}} \delta_{cc'} \theta_{x|c} (\delta_{xx'} - \theta_{x'|c}).
\end{split}\end{equation*}
Assume that $N_{+|c} \to \infty$ for every $c \in \C{C}$.
Because
$$N_{x|c} = \sum_{n=1 \atop c_n=c}^N \I_{x}(X_n),$$
where the $\I_{x}(X_n)$ can be seen as i.i.d.\ vectors in $\RN^{\C{X}}$, we can apply the multivariate
central limit theorem to conclude that
  \begin{equation}\label{eq:multinomial_clt}
    \sqrt{N_{+|c}} (\hat\theta_{x|c} - \theta_{x|c}) \leadsto N(0,\Sigma)
  \end{equation}
with covariance matrix $\Sigma$ with entries
  \begin{equation}\label{eq:multinomial_cov}
    (\Sigma)_{x|c,x'|c'} = \delta_{cc'} \theta_{x|c} (\delta_{xx'} - \theta_{x'|c}).
  \end{equation}

The central limit theorem thus provides the rate at which the ML estimate $\hat \theta$ converges to the true parameter $\theta$.
The likelihood ratio statistic $G_N$ is a function of $\hat \theta$ and $\hat\gamma$. 
The asymptotic analysis of $G_N$ can be obtained by performing a Taylor expansion of $g$ around the true $\theta$. 
This expansion can be done ``uniformly'' in $\hat\gamma$.

\begin{Lem}\label{lem:g_derivatives_H0}
  Let $\theta \in \Theta_0$ be positive, i.e., such that there exists a $\delta > 0$ with $\theta_{x|c} \ge \delta$ for all $x \in \C{X}, c \in \C{C}$.
  Let $\gamma \in \Theta_C$.
  For $\hat\theta$ with $\lVert \hat\theta - \theta \rVert$ sufficiently small,
  \begin{equation}\label{eq:g_expansion}
    g(\hat\theta, \gamma) = \frac{1}{2} (\hat\theta - \theta)^T \nabla^2_\theta g(\theta,\gamma) (\hat\theta - \theta) + o_P(\lVert\hat\theta - \theta\rVert^2),
  \end{equation}
  with
%  $g(\theta,\gamma) = 0$, 
%  $\nabla_\theta g(\theta,\gamma) = 0$, and
  $$\nabla^2_{\theta} g(\theta,\gamma) 
  %= \diag\bigg(\frac{1}{\theta_{X|\ast}}\bigg) \otimes \big(\diag(\gamma) - \gamma \gamma^T\big)
    = \diag\bigg(\frac{1}{\theta_{X|\ast}}\bigg) \otimes \Big(\diag(\sqrt{\gamma}) \big(I_{\C{C}} - \sqrt{\gamma} \sqrt{\gamma}^T\big)\diag(\sqrt{\gamma})\Big)$$
  where $\theta_{X|\ast} = (\theta_{x|c})_{x\in\C{X}}$ for an arbitrary $c \in \C{C}$.\footnote{Here, we used the Kronecker product notation for matrices, and $\diag(v)$ is a diagonal matrix with the components of vector $v$ on the diagonal.}
  The remainder term $o_P(\lvert\hat\theta - \theta\rVert^2)$ can be chosen to be a function $f(\theta,\hat\theta)$ that does not depend on $\gamma$, and for which
  $$\frac{f(\theta,\hat\theta)}{\lVert\theta - \hat\theta\rVert^2} \stackrel{P}{\to} 0.$$
\end{Lem}
\begin{proof}
The second order Taylor expansion of $g$ around the true $\theta$ is:
$$g(\theta + \epsilon, \gamma) = g(\theta,\gamma) + \epsilon^T \nabla_\theta g(\theta,\gamma) + \frac{1}{2} \epsilon^T \nabla^2_\theta g(\theta,\gamma) \epsilon + o(\lVert\epsilon\rVert^2),$$
where $\nabla_\theta g$ is the (partial) gradient of $g$ with respect to $\theta$, and
$\nabla_\theta^2 g$ is the Hessian of $g$ with respect to $\theta$, and the remainder
term $o(\lVert\epsilon\rVert^2)$ can be taken of the form $M \lVert\epsilon\rVert^3$ if the 
third-order partial derivatives of $g$ at $(\theta,\gamma)$ are bounded.
For a random $\epsilon = \hat\theta - \theta$ we obtain
  \begin{equation}\label{eq:g_expansion}
    g(\hat\theta, \gamma) - g(\theta,\gamma) = (\hat\theta - \theta)^T \nabla_\theta g(\theta,\gamma) + \frac{1}{2} (\hat\theta - \theta)^T \nabla^2_\theta g(\theta,\gamma) (\hat\theta - \theta) + o_P(\lVert\hat\theta - \theta\rVert^2),
  \end{equation}
where the remainder term is now random, and converges 
in probability to 0 at rate $\lVert\hat\theta - \theta\rVert^2$.
We will proceed by calculating the terms in the Taylor expansion.

The Kullback-Leibler divergence has the important property that $KL(P \,\Vert\, Q) \ge 0$ and $KL(P \,\Vert\, Q) = 0 \iff P = Q$ for all $P,Q \in \C{P}(\C{X})$. 
Together with the definition \eref{eq:G_test_statistics_limit}, 
this immediately implies that under $H_0$, $g(\theta,\gamma) = 0$ for all $\gamma \in \Theta_C$.

The gradient $\nabla_\theta g$ of $g$ w.r.t.\ $\theta$ has components:
  $$\frac{\partial g}{\partial \theta_{x|c}} 
%  = \gamma_c \log \frac{\theta_{x|c}}{\sum_{c'\in\C{C}} \gamma_c' \theta_{x|c'}} + \gamma_c - \sum_{c'\in\C{C}} \gamma_{c'} \theta_{x|c'} \frac{\gamma_c}{\sum_{c''\in\C{C}} \gamma_{c''} \theta_{x|c''}}
  = \gamma_c \log \frac{\theta_{x|c}}{\sum_{c'\in\C{C}} \gamma_c' \theta_{x|c'}}.$$
Under $H_0$, $\theta_{x|c} = \theta_{x|\ast}$ for all $x\in\C{X},c\in\C{C}$, and it follows that
the gradient vanishes for all $\gamma\in\Theta_C$.

Next, the Hessian w.r.t.\ $\theta$:
  $$\frac{\partial^2 g}{\partial \theta_{x|c} \partial \theta_{x'|c'}} = 
  \gamma_c \left( \frac{1}{\theta_{x|c}} \delta_{xx'} \delta_{cc'} - \frac{\gamma_{c'}}{\sum_{c''\in\C{C}} \gamma_{c''} \theta_{x|c''}} \delta_{xx'} \right).
  $$
Under $H_0$, $\theta_{x|c} = \theta_{x|\ast}$ for all $x\in\C{X},c\in\C{C}$, this simplifies to
  $$\frac{\partial^2 g}{\partial \theta_{x|c} \partial \theta_{x'|c'}} = 
  \frac{1}{\theta_{x|\ast}} \gamma_c \left( \delta_{xx'} \delta_{cc'} - \gamma_{c'} \delta_{xx'} \right).
  $$
By using the Kronecker product notation, this can be written as stated in the lemma.

Finally, to obtain the remainder term, we calculate the third order partial derivatives:
  $$\frac{\partial^3 g}{\partial \theta_{x|c} \partial \theta_{x'|c'} \partial \theta_{x''|c''}} = 
  \gamma_c \delta_{x,x''} \delta_{x,x'} \left( -\frac{1}{\theta_{x|c}^2} \delta_{c,c''} \delta_{cc'} + \frac{\gamma_{c'}\gamma_{c''} }{(\sum_{c'''\in\C{C}} \gamma_{c'''} \theta_{x|c'''})^2} \right).
  $$
These can be bounded uniformly in $\gamma$, using the assumption that all components of $\theta$ are bounded away from zero.
\end{proof}
We are now ready to prove the following result on the asymptotic distribution of $G_N$ under the null hypothesis, which is (surprisingly?) similar to Proposition~\ref{prp:G_test_asymptotics_H0}.
\begin{Prp}\label{prp:ext_G_test_asymptotics_H0}
  Let $\theta \in \Theta$ be positive, i.e., such that $\theta_{x|c} > 0$ for all $x \in \C{X}, c \in \C{C}$.
  Assume that $N_{+|c} \to \infty$ for all $c \in \C{C}$.
  Under $H_0 : X \Indep_{K(X \given C)} C$, the likelihood ratio test statistic \eref{eq:ext_G_test_statistic} converges
  to a $\chi^2$ distribution,
  $$G_N \leadsto \chi^2_\nu,$$
  with $\nu = (|\C{X}| - 1)(|\C{C}| - 1)$ degrees of freedom.
\end{Prp}
\begin{proof}
  Under $H_0$, Lemma \ref{lem:g_derivatives_H0} yields that:
%  $$g(\hat\theta, \hat\gamma) - g(\theta,\hat\gamma) = (\hat\theta - \theta)^T g'(\theta,\hat\gamma) + \frac{1}{2} (\hat\theta - \theta)^T g''(\theta,\hat\gamma) (\hat\theta - \theta) + o_P(\lVert\hat\theta - \theta\rVert^2)$$
%  i.e.,
  $$g(\hat\theta, \hat\gamma) = \frac{1}{2} (\hat\theta - \theta)^T \nabla_\theta^2 g(\theta,\hat\gamma) (\hat\theta - \theta) + o_P(\lVert|\hat\theta - \theta\rVert^2)$$
  where
  $$\nabla^2_{\theta} g(\theta,\hat\gamma) = \diag\bigg(\frac{1}{\theta_{X|\ast}}\bigg) \otimes \left(\diag(\sqrt{\hat\gamma}) \big(I_{\C{C}} - \sqrt{\hat\gamma} \sqrt{\hat\gamma}^T\big)\diag(\sqrt{\hat\gamma})\right)$$
  and the remainder term does not depend on $\hat\gamma$.
  Hence,
  \begin{equation*}\begin{split}
    & N (\hat\theta - \theta)^T \nabla_\theta^2(\theta,\hat\gamma) (\hat\theta - \theta) \\
    & = \sqrt{N} (\hat\theta - \theta)^T \diag\bigg(\frac{1}{\sqrt{\theta_{X|\ast}}} \otimes \sqrt{\hat\gamma}\bigg) \Big(I_{\C{X}} \otimes \big(I_{\C{C}} - \sqrt{\hat\gamma} \sqrt{\hat\gamma}^T\big)\Big) \diag\bigg(\frac{1}{\sqrt{\theta_{X|\ast}}} \otimes \sqrt{\hat\gamma}\bigg) (\hat\theta - \theta) \sqrt{N} \\
    & = \bigg\lVert \Big(I_{\C{X}} \otimes \big(I_{\C{C}} - \sqrt{\hat\gamma} \sqrt{\hat\gamma}^T\big)\Big) \diag\bigg(\frac{1}{\sqrt{\theta_{X|\ast}}} \otimes \sqrt{\hat\gamma}\bigg) (\hat\theta - \theta) \sqrt{N} \bigg\rVert^2.
  \end{split}\end{equation*}
%  Let $U \in SO(\RN^{\C{X}})$ be a rotation that brings $\sqrt{\theta_{X|\ast}}$ to $e_1$, and 
%  $\hat V \in SO(\RN^{\C{C}})$ a rotation that brings $\sqrt{\hat\gamma}$ to $e_1$. Then
%  $U \otimes \hat V \in SO(\RN^{\C{X}\times\C{C}})$, and hence preserves the Euclidean norm in $\RN^{\C{X}\times\C{C}}$.
%  Therefore,
%  \begin{equation*}\begin{split}
%    & \lVert \Big(I_{\C{X}} \otimes \big(I_{\C{C}} - \sqrt{\hat\gamma} \sqrt{\hat\gamma}^T\big)\Big) \diag(\frac{1}{\sqrt{\theta_{X|\ast}}} \otimes \sqrt{\hat\gamma}) (\hat\theta - \theta) \sqrt{N} \rVert^2 \\
%    & = \lVert (U \otimes \hat V) \Big(I_{\C{X}} \otimes \big(I_{\C{C}} - \sqrt{\hat\gamma} \sqrt{\hat\gamma}^T\big)\Big) (U^T \otimes I) (U \otimes I)\diag(\frac{1}{\sqrt{\theta_{X|\ast}}} \otimes \sqrt{\hat\gamma}) (\hat\theta - \theta) \sqrt{N} \rVert^2 \\
% \end{split}\end{equation*}

  Under $H_0$, \eref{eq:multinomial_clt} can be written as
  $$\sqrt{N} \diag(1_{\C{X}} \otimes \sqrt{\hat\gamma}) (\hat\theta - \theta)  \leadsto N(0,\Sigma),$$
  where $\otimes$ denotes the Kronecker product of two vectors, in this case the all-ones vector $1_{\C{X}} \in \RN^{\C{X}}$
  and the vector $\sqrt{\hat\gamma} \in \RN^{\C{C}}$.
  Furthermore, under $H_0$, $\theta = \theta_{X|\ast} \otimes 1_{\C{C}}$, and
  the covariance matrix $\Sigma$ given in \eref{eq:multinomial_cov} simplifies to
    %$$(\Sigma)_{x|c,x'|c'} = \delta_{cc'} \theta_{x|c} (\delta_{xx'} - \theta_{x'|c}).$$
  $$\Sigma = \big(\theta_{X|\ast} (I_{\C{X}} - \theta_{X|\ast}^T)\big) \otimes I_{\C{C}}
           = \Bigg(\diag\bigg(\sqrt{\theta_{X|\ast}}\bigg) \bigg(I_{\C{X}} - \sqrt{\theta_{X|\ast}}\sqrt{\theta_{X|\ast}^T}\bigg) \diag\bigg(\sqrt{\theta_{X|\ast}}\bigg)\Bigg) \otimes I_{\C{C}}.$$
  Hence, by rescaling with $\sqrt{\theta}$,
  $$\sqrt{N} \diag\bigg(\frac{1}{\sqrt{\theta_{X|\ast}}} \otimes \sqrt{\hat\gamma}\bigg) (\hat\theta - \theta) \leadsto N(0,\tilde\Sigma),$$
  with
  $$\tilde\Sigma = \Big(I_{\C{X}} - \sqrt{\theta_{X|\ast}} \sqrt{\theta_{X|\ast}^T}\Big) \otimes I_{\C{C}}.$$
  Note that $\tilde\Sigma$ is an orthogonal projection, hence $\tilde\Sigma = \tilde\Sigma^T \tilde\Sigma$.
  Also
  $$\hat\Gamma := I_{\C{X}} \otimes \Big(I_{\C{C}} - \sqrt{\hat\gamma} \sqrt{\hat\gamma}^T\Big)$$
  is an orthogonal projection.

  We now use the fact that if $S_N \leadsto N(0,P \otimes I)$ for an orthogonal projection $P$ on a $\nu$-dimensional
  subspace, and $Q_N$ are orthogonal projections on $\mu$-dimensional subspaces, then
  $\lVert (Q_N \otimes I) S_N \rVert \leadsto \chi^2_{\mu\nu}$. \Joris{Elaborate?}
  
  In the case at hand, we note that while $\hat\gamma$ yields a rotation in $\RN^{\C{C}}$ that could be different for
  every $N$, this does not affect the limiting distribution, because of the rotational symmetry of the multivariate
  distribution. Note here that the Kronecker product interacts nicely with rotations and projections. The influence
  of $\hat\gamma$ and $\theta_{X|\ast}$ are decoupled by the Kronecker product structure.
  We conclude that
    $$N (\hat\theta - \theta)^T \nabla_\theta^2(\theta,\hat\gamma) (\hat\theta - \theta) \leadsto \chi^2_{\nu}$$
  with $\nu = (|\C{X}|-1)(|\C{C}|-1)$.
  With Slutsky's Lemma, also
  $$G_N = 2N g(\hat\theta,\hat\gamma) = N (\hat\theta - \theta)^T \nabla_\theta^2(\theta,\hat\gamma) (\hat\theta - \theta) + o_P(\lVert|\hat\theta - \theta\rVert^2) \leadsto \chi^2_{\nu}.$$
\end{proof}
So, even the asymptotic distribution under the null hypothesis is the same as for the case of two random variables, 
even though we used a different parameterization,
and we relaxed the assumption that the $c_n$'s are i.i.d.:
we only assumed that $N_{+|c} \to \infty$ for each $c$.\footnote{If some of the $N_{+|c}$ stay finite asymptotically, then these $c$'s can be ignored, and we still get asymptotically a chi-square distribution, but with less degrees of freedom.}
In particular, this result applies also to the case when
$C$ is a random variable. Hence, we have reobtained Proposition~\ref{prp:G_test_asymptotics_H0} as a special case.

So, summarizing, we started from quite a different sampling scheme, and did not treat $C$ as a random variable, yet we
ended up with exactly the same likelihood ratio test for testing $X \Indep_{K(X|C)} C$ as we derived 
for testing the independence $X \Indep_{P(X,Y)} Y$ between two random variables.

What about the consistency of this test? To show consistency, it turns out that we need a stronger assumption
than $N_{+|c} \to \infty$, but it will still be weaker than the i.i.d.\ assumption we made for the case that $C$ is random.
\begin{Cor}\label{cor:ext_G_test_consistency}
Consider an infinite sequence of $G$ tests performed on the first $N$ samples of an infinitely large
data set $(X_n,c_n)_{n=1}^\infty$,
where one accepts $H_1 : X \nIndep C$ if $G_N \ge \tau_N$, and otherwise accepts $H_0 : X \Indep C$, 
for some given sequence of thresholds $\tau_N$.
Assume that $\theta \in \Theta$ is positive, i.e., such that $\theta_{x|c} > 0$ for all $x \in \C{X}, c \in \C{C}$.
Assume further that the fractions $N_{+|c} / N \to \infty$ are bounded away from zero asymptotically,
i.e., there exists $\epsilon > 0$ such that for all $c \in \C{C}$, $N_{+|c} / N > \epsilon$ for large $N$.
Then this sequence of tests is asymptotically consistent if $\tau_N \to \infty$ but $\tau_N / N \to 0$.
\end{Cor}
\begin{proof}
Under $H_1$, there is a pair $c,c' \in \C{C}$ such that $\theta_{X|c} \ne \theta_{X|c'}$, and hence,
  $$\delta = KL\Big(\theta_{X|c} \,\Vert\, \sum_{c\in\C{C}} \gamma_c \theta_{X|c}\Big) > 0.$$
Therefore, $g(\theta,\hat\gamma) \ge \epsilon \delta > 0$ for large $N$.
If $\hat\theta \stackrel{a.s.}{\to} \theta$, then by the uniform \Joris{check!} continuity of $g$, 
also $g(\hat\theta,\hat\gamma)$ a.s.\ has a positive lower bound for large $N$.
The same then holds for $G_N / N = 2g(\hat\theta,\hat\gamma)$.

To see that $\hat\theta \stackrel{a.s.}{\to} \theta$, we can apply the strong law of large numbers.
Since the $X_n$ are assumed to be conditionally i.i.d.\ given $c_n$, and 
  $$\Exp \big(\I_{(x,c)}(X_n,c_n)\big) = \theta_{x|c} \I_c(c_n)$$
for all $x \in \C{X}, c \in \C{C}$, and $N_{+|c} \to \infty$ for all
$c \in \C{C}$, we conclude that $N_{x|c} / N_{+|c} \stackrel{a.s.}{\to} \theta_{x|c}$ 
for all $x \in \C{X}, c \in \C{C}$
by the strong law of large numbers. Hence, $\hat\theta \stackrel{a.s.}{\to} \theta$.

We can now reason analogously as in the proof of 
Corollary~\ref{cor:G_test_consistency} to conclude that
the probability of a Type II error vanishes asymptotically.

For an asymptotic estimate of the probability of a Type I error,
we can make use of Proposition~\ref{prp:ext_G_test_asymptotics_H0},
which states that $G_N \leadsto \chi_{\nu}^2$. 
This part of the proof is identical to the corresponding part of the proof of Corollary~\ref{cor:G_test_consistency}.
\end{proof}
The conditions in this corollary are sufficient, but not necessary. For example,
not all the rates $N_{+|c}$ have to be lower bounded, it suffices if this is
the case for a subset of $\C{C}$ for which the distributions $K(X \mid C=c)$ differ.
It also shows how consistency could fail: e.g., if the distributions
$K(X \mid C=c)$ only differ on some subset of $\C{C}$, but that subset is not 
observed sufficiently often asymptotically. This is in line with the intuition that
when testing for the presence of a causal effect of $C$ on $X$ in a controlled
setting (not necessarily randomized, i.e., as in Proposition~\ref{prp:RCT1}), if
nothing is known about how $X$ might depend on $C$, it is best to gather sufficient
data for each value that $C$ can take.

%\Joris{For consistency of this test, we need that $N_{+|c} \to \infty$ for each value of $c$.
%In the binary case, suppose $N_{+|0} \to\infty$ and $N_{+|1}$ stays fixed, then we obtain
%a likelihood ratio test statistic for the test $H_0 : p_{+|1} = q_0$ based only on the data for $c=1$,
%where $q_0 := \lim_{N\to\infty} N_{+|0} / N$. That is nice.
%}


\subsubsection{The general categorical case}

Finally, let us consider the most general case of testing a conditional independence involving both random and
non-random variables. Again, we will restrict ourselves to the case that all variables take values in finite spaces.
We will formulate a general version of the $G$ test.

%Let $\langle \C{W} \times \C{T}, K(W\mid T) \rangle$ be a finite transition space (i.e., with $|\C{W}| < \infty$ and
%$|\C{T}| < \infty$) with Markov kernel $K(W \mid T): \C{T} \dshto \C{W}$. Consider random fields $X,Y,Z,D$ with common
%domain $\C{W} \times \C{T}$ and (finite) codomains $\C{X},\C{Y},\C{Z}$, respectively, and $T : \C{W} \times \C{T} \to \C{T}$
%the canonical projection map. 

Suppose we have five variables: three random variables $X,Y,Z$ and two 
non-random input variables $C,D$, taking values in spaces $\C{X}, \C{Y}, \C{Z}, \C{C}, \C{D}$, respectively. 
Assume that all the spaces $\C{X},\C{Y},\C{Z},\C{C},\C{D}$ are finite.
Suppose we have a kernel $K(X,Y,Z,\dots|C,D) : \C{C} \times \C{D} \dshto \C{X} \times \C{Y} \times \C{Z}$.
We formulate a statistical test for testing the conditional independence
$$H_0 : X \Indep_{K(X,Y,Z,\dots|C,D)} Y, C \given Z, D$$
against the alternative
$$H_1 : X \nIndep_{K(X,Y,Z,\dots|C,D)} Y, C \given Z, D.$$
The null hypothesis is equivalent, by definition, to the existence of a kernel $K(X|Z,D)$ such that
$$K(X,Y,Z \mid C,D) = K(X \mid Z,D) \otimes K(Y,Z \mid C,D).$$
%In terms of transition probabilities:
%$$\tilde{K}_{x,y,z|c,d} = \tilde{L}_{x|z,d} \tilde{K}_{+,y,z|c,d}$$
We will not work out the details, as these are analogous to what we have seen before, but will directly formulate
the likelihood ratio test statistic:
\begin{equation}\label{eq:general_G_test_statistic}
  G_N = 2 \sum_{x\in\C{X}}\sum_{y\in\C{Y}}\sum_{z\in\C{Z}}\sum_{c\in\C{C}}\sum_{d\in\C{D}}
  N_{xyzcd} \log \frac{N_{xyzcd} N_{++z+d}}{N_{x+z+d} N_{+yzcd}}.
\end{equation}
With the correspondence $(X,(Y,C),(Z,D)) \leftrightarrow (X,Y,Z)$, this likelihood
ratio statistics is seen to be identical to \eref{eq:cond_G_test_statistic}, the one for the case
of three random variables. Note that this likelihood ratio treats $C$ and $Y$ on equal footing, as
well as $D$ and $Z$.  This, again, suggests that for the asymptotic analysis we will obtain
a similar result as that for the conditional $G$ test with purely random variables, albeit under milder
assumptions on the sampling scheme for $C$ and $D$. We will not work out the details here, but directly
formulate the result.

Parameterize the kernel $K(X,Y,Z \mid C,D)$ as
$$K(X=x,Y=y,Z=z \mid C=c,D=d) = \theta_{xyz|cd}$$
with $\theta$ in 
$$\Theta = \{ \theta \in \prod_{(x,y,z,c,d) \in \C{X}\times\C{Y}\times\C{Z}\times\C{C}\times\C{D}} : \theta_{+++|cd} = 1 \ \forall c\in\C{C}, d\in\C{D} \}.$$
\begin{Prp}\label{prp:general_G_test_asymptotics_H0}
  Let $\theta \in \Theta$ be positive, i.e., such that $\theta_{xyz|cd} > 0$ for all $x \in \C{X}, y \in \C{Y}, z \in \C{Z}, c \in \C{C}, d \in \C{D}$.
  Assume that $N_{+++|cd} \to \infty$ for all $c \in \C{C}, d \in \C{D}$.
  Under $H_0 : X \Indep_{K(X,Y,Z \given C,D)} Y,C \given Z,D$, the likelihood ratio test statistic \eref{eq:general_G_test_statistic} converges
  to a $\chi^2$ distribution,
  $$G_N \leadsto \chi^2_\nu,$$
  with $\nu = |\C{Z}|\,|\C{D}|\,(|\C{X}| - 1)(|\C{Y}||\C{C}| - 1)$ degrees of freedom.
\end{Prp}
\begin{proof}
  Note that the likelihood ratio test statistic \eref{eq:general_G_test_statistic} is a sum
  over $z \in \C{Z}$ and $d \in \C{D}$ of a likelihood ratio test statistic of the form 
  \eref{eq:ext_G_test_statistic} (where $(Y,C)$ in the former corresponds with $C$ in the latter). 
  \Joris{Elaborate?}
\end{proof}
%It is possible to drop the positivity assumption for $\theta$. This leads to a reduction in the degrees of
%freedom, but the statement that $G_N$ converges in distribution to a chi-square distribution remains valid,
%but the number of degrees of freedom of that distribution is reduced.
We also obtain a similar result as before on the asymptotic consistency.
\begin{Cor}\label{cor:gen_G_test_consistency}
Consider an infinite sequence of $G$ tests performed on the first $N$ samples of an infinitely large
data set $(X_n,Y_n,Z_n,c_n,d_n)_{n=1}^\infty$,
where one decides
  $$\begin{cases}
    H_0 : X \Indep_{K(X,Y,Z \given C,D)} Y,C \given Z,D     & \text{if $G_N < \tau_N$},\\
    H_1 : X \nIndep_{K(X,Y,Z \given C,D)} Y,C \given Z,D & \text{if $G_N \ge \tau_N$,}
  \end{cases}$$
for some given sequence of thresholds $\tau_N$.
Assume that $\theta \in \Theta$ is positive, i.e., such that $\theta_{xyz|cd} > 0$ for all $x \in \C{X}, y \in \C{Y}, z \in \C{Z}, c \in \C{C}, d \in \C{D}$.
Assume further that the fractions $N_{+++|cd} / N \to \infty$ are bounded away from zero asymptotically,
i.e., there exists $\epsilon > 0$ such that for all $c \in \C{C}, d \in \C{D}$, $N_{+++|cd} / N > \epsilon$ for large $N$.
Then this sequence of tests is asymptotically consistent if $\tau_N \to \infty$ but $\tau_N / N \to 0$.
\end{Cor}

%CI testing: $X \Indep Y, J_1 \mid Z, J_2$ (with $J = J_1 \dot\cup J_2$)
%    iff $\forall j_2: X \Indep Y, J_1 \mid Z, j_2$
%    iff $\forall j_2: P(X \mid Y,J_1,Z,j_2) = P(X \mid Z, j_2)$

\begin{comment}
\subsubsection{Leftovers}

Miller and Madow (1954) give the asymptotic distributions for the plug-in estimator of Shannon entropy.
While for generic parameter choices, this just yields the expected normal distribution, funny stuff
happens if some of the $p_x$ are zero and at the same time, the non-zero ones are uniform. Then one
gets a chi-squared distribution instead. 
%@report{MillerMadow1954,
%  author = {G. A. Miller and W. G. Madow},
%  title = {On the maximum likelihood estimate of the {S}hannon-{W}eaver measure of information},
%  institution = {Air Force Cambridge Research Center},
%  number = {54-75},
%  year = 1954
%}

\url{https://stats.stackexchange.com/questions/4816/what-are-the-sharpest-known-tail-bounds-for-chi-k2-distributed-variables}
mentions a bound 
$$P(\chi_{\nu}^2 \ge \nu + 2\sqrt{\nu x} + 2x) \le e^{-x}$$
which I quickly turned into
$$\chi^2_{\nu,1-\alpha} \le \nu + 2\sqrt{\nu \log \alpha^{-1}} + 2 \log \alpha^{-1}$$
This should allow us to find a sufficient condition for $\chi^2_{\nu,1-\alpha_N} / N \to 0$.
Not sure if this is entirely correct.

\begin{Cor}
Consider the sequence of $G$-tests where one rejects $H_0 : X \Indep Y$ if $G_N \ge \tau_N$ for some given
sequence of thresholds $\tau_N$.
Under the regularity assumption \eref{eq:G_test_regularity}, this sequence is
asymptotically valid with level $\alpha$ if $\tau_N \to \chi_{\nu,1-\alpha}^2$,
where $\nu = (|\C{X}|-1)(|\C{Y}|-1)$.
\end{Cor}
\begin{proof}
By Proposition~\ref{prp:G_test_asymptotics_H0}, under $H_0$,
$G_N \leadsto \chi_{\nu}^2$. Since the distribution function of $\chi_{\nu}^2$ is continuous,
this implies uniform convergence of the distribution functions:
  $$\sup_{x \in \RN} |F_{G_N}(x) - F_{\chi_{\nu}^2}(x)| \to 0.$$
Hence
  $$|F_{G_N}(\tau_N) - F_{\chi_{\nu}^2}(\tau_N)| \le \sup_{x\in\RN} |F_{G_N}(x) - F_{\chi_{\nu}^2}(x)| \to 0.$$
Since $\tau_N \to F^{-1}_{\chi_{\nu}^2}(1-\alpha)$ and $F_{\chi_{\nu}^2}$ is continuous,
  $$F_{\chi_\nu^2}(\tau_N) \to F_{\chi_\nu^2}(F^{-1}_{\chi_\nu^2}(1-\alpha)) = 1-\alpha.$$
Combining these two, we conclude
  $$F_{G_N}(\tau_N) \to 1-\alpha$$
and hence the probability of a Type I error converges to $\alpha$,
  $$P(G_N \ge \tau_N) \to \alpha.$$
\end{proof}

To obtain the asymptotic behavior of the distribution of the test statistic under the alternative hypothesis,
one can find a general analysis in \cite{LomnickiZaremba1959} (for arbitrary $\theta_{xy} \in \Theta$). 
Here we will again make the
simplifying assumption that $\theta_{xy} > 0$ for all $x\in\C{X},y\in\C{Y}$.
\begin{Prp}\label{prp:G_test_asymptotics_H1}
Under $H_1 : X \nIndep Y$ (equivalently, if $I := I(X;Y) > 0$), and with regularity assumption \eref{eq:G_test_regularity},
  $$\sqrt{N}(\hat I_N - I) \leadsto N(0,\sigma^2)$$
where
  $$\sigma^2 := \sum_{x\in\C{X}} \sum_{y\in\C{Y}} \theta_{xy} \left(\log\frac{\theta_{xy}}{\theta_{x+} \theta_{+ y}}\right)^2 - I^2$$
and $I = I(X;Y)$ as sample size $N \to \infty$.
\end{Prp}
\begin{proof}
%  \begin{equation*}\begin{split}
%    &\Cov (\I_{xy}(X_n,Y_n), \I_{x'y'}(X_n,Y_n)) \\
%    &= \Exp(\I_{xy}(X_n,Y_n)\I_{x'y'}(X_n,Y_n)) - \Exp(\I_{xy}(X_n,Y_n)) \Exp(\I_{x'y'}(X_n,Y_n))\\
%    &= \delta_{xy,x'y'} \theta_{xy} - \theta_{xy} \theta_{x'y'} \\
%    &=: \Sigma
%  \end{split}\end{equation*}
  For $x,x' \in \C{X}$, $y,y' \in \C{Y}$, we calculate:
  \begin{equation*}\begin{split}
    &N \Cov (\hat\theta_{xy}, \hat\theta_{x'y'}) \\
    &= N \Cov \left(\frac{1}{N} \sum_{n=1}^N \I_{xy}(X_n,Y_n), \frac{1}{N} \sum_{n'=1}^N \I_{x'y'}(X_{n'},Y_{n'})\right) \\
    &= \frac{1}{N} \sum_{n=1}^N \sum_{n'=1}^N \Cov (\I_{xy}(X_n,Y_n), \I_{x'y'}(X_{n'},Y_{n'})) \\
    &= \Cov (\I_{xy}(X_1,Y_1), \I_{x'y'}(X_{1},Y_{1})) \\
    &= \theta_{xy} (\delta_{xy,x'y'} - \theta_{x'y'}) \\
    &=: \Sigma_{xy,x'y'}
  \end{split}\end{equation*}
  Asymptotically, therefore,
  $$\sqrt{N}(\hat\theta - \theta) \leadsto N(0,\Sigma)$$
  with the multivariate central limit theorem.

  With the multivariate delta method we can obtain the limiting distribution of the sequence
  $\sqrt{N}(I(\hat\theta) - I(\theta))$. The delta method consists of an asymptotic linearization
  around $\theta$, and asserts that the limiting distribution is the same as that of
  $\sqrt{N}I'(\theta)(\hat\theta - \theta)$. In this case, we obtain
  $$\sqrt{N}I'(\theta)(\hat\theta - \theta) \leadsto %I'(\theta) N(0,\Sigma) = 
  N\big(0,I'(\theta)^T \Sigma I'(\theta)\big).$$
  The gradient of
  $$\theta \mapsto I(\theta) := \sum_{x\in\C{X}}\sum_{y\in\C{Y}} \theta_{xy} \log \frac{\theta_{xy}}{\theta_{x +} \theta_{+ y}}
       = \sum_{x\in\C{X}}\sum_{y\in\C{Y}} \theta_{xy} \log \theta_{xy}
       - \sum_{x\in\C{X}} \theta_{x+} \log \theta_{x +}
       - \sum_{y\in\C{Y}} \theta_{+ y} \log \theta_{+ y}$$
  is given by
%  = \sum_{x\in\C{X}}\sum_{y\in\C{Y}} \theta_{xy} \log \theta_{xy}
%  - \sum_{x\in\C{X}}\sum_{y\in\C{Y}} \theta_{xy} \log \sum_{y'} \theta_{xy'} 
%  - \sum_{x\in\C{X}}\sum_{y\in\C{Y}} \theta_{xy} \log \sum_{x'} \theta_{x'y}$$
  $$I'(\theta) = \frac{\partial I(\theta)}{\partial \theta_{xy}} = 
    \log \theta_{xy} + 1
    - \log \theta_{x +} - 1
    - \log \theta_{+ y} - 1
   = \log \frac{\theta_{xy}}{\theta_{x +} \theta_{+ y}} - 1$$
  Hence
   \begin{equation*}\begin{split}
   & I'(\theta)^T \Sigma I'(\theta) \\
   & = \sum_{xy} \sum_{x'y'} \left( \log \frac{\theta_{xy}}{\theta_{x +} \theta_{+ y}} - 1 \right)
     \Sigma_{xy,x'y'}
     \left( \log \frac{\theta_{x'y'}}{\theta_{x' +} \theta_{+ y'}} - 1 \right)\\
   & = \sum_{xy} \sum_{x'y'} \left( \log \frac{\theta_{xy}}{\theta_{x +} \theta_{+ y}} - 1 \right)
     (\delta_{xy,x'y'} \theta_{xy} - \theta_{xy} \theta_{x'y'})
     \left( \log \frac{\theta_{x'y'}}{\theta_{x' +} \theta_{+ y'}} - 1 \right)\\
   & = \sum_{xy} \theta_{xy} \left( \log \frac{\theta_{xy}}{\theta_{x +} \theta_{+ y}} - 1 \right)^2 -
     \sum_{xy} \theta_{xy} \left( \log \frac{\theta_{xy}}{\theta_{x +} \theta_{+ y}} - 1 \right) \sum_{x'y'} 
     \theta_{x'y'}
     \left( \log \frac{\theta_{x'y'}}{\theta_{x' +} \theta_{+ y'}} - 1 \right) \\
   & = \sigma^2.
   \end{split}\end{equation*} 
   which completes the proof.
\end{proof}
Note that this is not very useful in practice to estimate the Type II error if one does not know the true $I(\theta)$.

We provide a more direct proof of Proposition~\ref{prp:G_test_asymptotics_H0} as well, 
in the hope that it will be useful in generalizing
the analysis later on. It is based on the delta method, which exploits a Taylor expansion.
We start with a little lemma on the Taylor expansion of the mutual information.
In the rest of this subsection, we let $\theta \in \RN^{\C{X}\times \C{Y}}$ denote a vector
with components $\theta_{xy}$ for $x \in \C{X}$, $y \in \C{Y}$. We let $\theta_X \in \RN^{\C{X}}$
denote the vector with components $(\theta_X)_x = \theta_{x+} = \sum_{y\in\C{Y}} \theta_{xy}$, 
and $\theta_Y \in \RN^{\C{Y}}$ the vector with components $(\theta_Y)_y = \theta_{+ y} = \sum_{x\in\C{X}} \theta_{xy}$.
\begin{Lem}
  The gradient of the mutual information
  $$\theta \mapsto I(\theta) := \sum_{x\in\C{X}}\sum_{y\in\C{Y}} \theta_{xy} \log \frac{\theta_{xy}}{\theta_{x +} \theta_{+ y}}$$
%       = \sum_{x\in\C{X}}\sum_{y\in\C{Y}} \theta_{xy} \log \theta_{xy}
%       - \sum_{x\in\C{X}} \theta_{x+} \log \theta_{x +}
%       - \sum_{y\in\C{Y}} \theta_{+ y} \log \theta_{+ y}$$
  is given by
%  = \sum_{x\in\C{X}}\sum_{y\in\C{Y}} \theta_{xy} \log \theta_{xy}
%  - \sum_{x\in\C{X}}\sum_{y\in\C{Y}} \theta_{xy} \log \sum_{y'} \theta_{xy'} 
%  - \sum_{x\in\C{X}}\sum_{y\in\C{Y}} \theta_{xy} \log \sum_{x'} \theta_{x'y}$$
  $$\big(I'(\theta))_{xy} = \frac{\partial I(\theta)}{\partial \theta_{xy}}
%    \log \theta_{xy} + 1
%    - \log \theta_{x +} - 1
%    - \log \theta_{+ y} - 1
   = \log \frac{\theta_{xy}}{\theta_{x +} \theta_{+ y}} - 1$$
  In other words,
  $$I'(\theta) = \log \frac{\theta}{\theta_X \otimes \theta_Y} - 1.$$
  When applied on differences of the form $(\hat\theta - \theta) \in 1^\perp$, we get
  $$I'(\theta)(\hat\theta - \theta) = (\log \frac{\theta}{\theta_X \otimes \theta_Y})^T (\hat\theta - \theta).$$
  The Hessian is given by
   $$\big(I''(\theta)\big)_{xy,x'y'} = \frac{\partial^2 I(\theta)}{\partial \theta_{xy} \partial \theta_{x'y'}}
     = \frac{\partial}{\partial \theta_{x'y'}} \log \frac{\theta_{xy}}{\theta_{x +} \theta_{+ y}}
     = \frac{1}{\theta_{xy}} \delta_{xy,x'y'} - \delta_{x,x'} \frac{1}{\theta_{x +}} - \delta_{y,y'} \frac{1}{\theta_{+ y}}$$
\end{Lem}
\begin{proof}
  Direct calculation.
\end{proof}
We also make use of the following lemma regarding the categorical distribution.
\begin{Lem}
  Let $X \sim \mathrm{Cat}(\theta_1,\dots,\theta_k)$ have a categorical distribution with parameter
  $\theta \in \RN^k$ such that $\sum_{i=1}^k \theta_i = 1$. In other words, $P(X = x) = \theta_x$ for $x=1,\dots,k$. Then
  $$\Exp \big(\I_x(X)\big) = \theta_x$$
  and
  $$\Cov (\I_x(X), \I_{x'}(X)) = \theta_x (\delta_{x,x'} - \theta_{x'}).$$
\end{Lem}

Now we continue with the direct proof of Proposition~\ref{prp:G_test_asymptotics_H0}.
\begin{proof}
  For $x,x' \in \C{X}$, $y,y' \in \C{Y}$, we calculate:
  \begin{equation*}\begin{split}
    &N \Cov (\hat\theta_{xy}, \hat\theta_{x'y'}) \\
    &= N \Cov \left(\frac{1}{N} \sum_{n=1}^N \I_{xy}(X_n,Y_n), \frac{1}{N} \sum_{n'=1}^N \I_{x'y'}(X_{n'},Y_{n'})\right) \\
    &= \frac{1}{N} \sum_{n=1}^N \sum_{n'=1}^N \Cov (\I_{xy}(X_n,Y_n), \I_{x'y'}(X_{n'},Y_{n'})) \\
    &= \theta_{xy} (\delta_{xy,x'y'} - \theta_{x'y'}) \\
    &=: \Sigma_{xy,x'y'}
  \end{split}\end{equation*}
  Asymptotically, therefore,
  \begin{equation}\label{eq:asymptotics_categorial_ML_estimators}
    \sqrt{N}(\hat\theta - \theta) \leadsto N(0,\Sigma)
  \end{equation}
  with the multivariate central limit theorem.
%  with the help of e.g.\ Theorem 5.9 and the remarks in paragraph 5.4.2 in \cite{Bijma++2017}
%  (and \cite{Vandervaart1998} for the complete proof),
%  where again we have to choose a parameterization of the right dimensionality when checking the
%  regularity conditions.

  With the multivariate delta method we can obtain the limiting distribution of the sequence
  $\sqrt{N}(I(\hat\theta) - I(\theta))$. The delta method consists of an asymptotic linearization
  around $\theta$, and asserts that the limiting distribution is the same as that of
  $\sqrt{N}I'(\theta)(\hat\theta - \theta)$. In this case, we obtain
  $$\sqrt{N}I'(\theta)(\hat\theta - \theta) \leadsto %I'(\theta) N(0,\Sigma) = 
  N\big(0,I'(\theta)^T \Sigma I'(\theta)\big).$$

  With the Lemma on the gradient of the mutual information,
  under $H_0$, $I'(\theta)(\hat\theta - \theta) = 0$.
  Therefore, in that case we conclude that
  $$\sqrt{N}(I(\hat\theta) - I(\theta)) \leadsto \delta_{0}.$$

  Under $H_0$, we go beyond a first order Taylor expansion (see also section 3.4. in \cite{Vandervaart1998}).
  By rescaling (a trick due to Pearson), we obtain from \eref{eq:asymptotics_categorial_ML_estimators} that
  $$\sqrt{N}\frac{\hat\theta - \theta}{\sqrt{\theta}} \leadsto N(0,\tilde\Sigma)$$
  with
  $$\tilde\Sigma = \frac{1}{\sqrt{\theta}} \Sigma \frac{1}{\sqrt{\theta}}^T = I - \sqrt{\theta} \sqrt{\theta}^T.$$
   This is the projection on the orthocomplement of the linear subspace spanned by the vector $\sqrt{\theta}$, i.e.,
   a subspace we denote as $\llbracket\sqrt{\theta}\rrbracket^T$.
   Hence $N \lVert \frac{\hat\theta - \theta}{\sqrt{\theta}} \rVert^2 \leadsto \chi^2_{|\C{X}||\C{Y}|-1}$
   (which, incidentally, provides the asymptotic validity analysis for the so-called $\chi^2$-test).
   Furthermore, since $\tilde\Sigma$ is a projection, also
   $$\sqrt{N} \tilde{\Sigma} \frac{\hat\theta - \theta}{\sqrt{\theta}} \leadsto N(0,\tilde\Sigma)$$
     
   Under $H_0$, $I'(\theta) = 0$ and hence we need to go beyond the first order expansion.
   $$I''(\theta) = \frac{\partial^2 I(\theta)}{\partial \theta_{xy} \partial \theta_{x'y'}}
     = \frac{\partial}{\partial \theta_{x'y'}} \log \frac{\theta_{xy}}{\theta_{x +} \theta_{+ y}}
     = \frac{1}{\theta_{xy}} \delta_{xy,x'y'} - \delta_{x,x'} \frac{1}{\theta_{x +}} - \delta_{y,y'} \frac{1}{\theta_{+ y}}$$
   Then, by the second-order delta method,
   $$N (I(\hat\theta) - I(\theta)) \leadsto \frac{1}{2} \sqrt{N}(\hat\theta - \theta)^T I''(\theta) \sqrt{N}(\hat\theta - \theta)$$
   With the rescaling trick, we can write:   
   $$N (I(\hat\theta) - I(\theta)) \leadsto \frac{1}{2} \sqrt{N}\frac{\hat\theta - \theta}{\sqrt{\theta}}^T \big(\sqrt{\theta}^T I''(\theta) \sqrt{\theta}\big) \sqrt{N} \frac{\hat\theta - \theta}{\sqrt{\theta}}$$
   Making use of the Kronecker product notation for matrices, 
   and denoting the $|\C{X}|$-dimensional vector with components $\theta_{x +}$ as $\theta_X$ and the
   $|\C{Y}|$-dimensional vector with components $\theta_{+ y}$ as $\theta_Y$, we can write
   $$\sqrt{\theta}^T I''(\theta) \sqrt{\theta} = I \otimes I - I \otimes (\sqrt{\theta_Y} \sqrt{\theta_Y}^T) - (\sqrt{\theta_X}\sqrt{\theta_X}^T) \otimes I.$$
   It has an eigenspace with eigenvalue 0 spanned by vectors of the form
   $\sqrt{\theta_X}^\perp \otimes \sqrt{\theta_Y}$ and
   $\sqrt{\theta_X} \otimes \sqrt{\theta_Y}^\perp$.
   It also has eigenvalue $-1$ corresponding to the space spanned by $\sqrt{\theta_X} \otimes \sqrt{\theta_Y}$.
   Finally, it has eigenvalue 1 corresponding to the space spanned by
   $\sqrt{\theta_X}^\perp \otimes \sqrt{\theta_Y}^\perp$.
   %and $\sqrt{\theta_X} \otimes \sqrt{\theta_Y}$.
   After ``sandwiching'' with $\tilde{\Sigma}$, we conclude that
   $$\tilde{\Sigma}^T \sqrt{\theta}^T I''(\theta) \sqrt{\theta} \tilde{\Sigma}$$
   has eigenvalue 1 with multiplicity $(|\C{X}|-1)(|\C{Y}|-1)$, and eigenvalue 0 with multiplicity $|\C{X}| + |\C{Y}| - 1$.
   Hence, 
   $$2N (I(\hat\theta) - I(\theta)) \leadsto \chi^2_{\nu},$$
   and since $I(\theta)=0$ under $H_0$,
   $$2N I(\hat\theta) \leadsto \chi^2_{\nu}.$$
\end{proof}

\Joris{Overview:
$$\hat\theta \stackrel{a.s.}{\to} \theta$$
$$I(\hat\theta) \stackrel{a.s.}{\to} I(\theta)$$
hence:
$$G_N/N \stackrel{a.s.}{\to} 2 I(\theta)$$
Delta method:
$$\sqrt{N}(I(\hat\theta)-I(\theta)) \leadsto \delta_0$$
hence under $H_0$:
$$G_N/\sqrt{N} \leadsto \delta_0$$
Second order delta method under $H_0$:
$$G_N \leadsto \chi^2_{\nu}.$$
}
\end{comment}

